\chapter{Composition of paths, records, and signatures}
\label{chap:composicion-especies-rev6}
\index{composition of paths}
\index{composition of records}
\index{meson!compositional expression}
\index{baryon!compositional expression}

A composite species cannot enter the law of masses as a name to which
a row of five integers is retrospectively assigned. Before the signature,
an object capable of composition must exist. The atlas of
\cref{chap:atlas-genealogia-masas} supplies the finite alphabet; the
surviving extension and register of
\cref{chap:extension-ruta-caracter} supply the memory that may
be composed. The question of this chapter is to construct the intermediate stratum
without consulting target masses:
\[
\begin{gathered}
 \text{typed compositional configuration}
 \longrightarrow
 \text{complete records of its constituents}\\
 \longrightarrow
 \text{composite record}
 \longrightarrow
 \ZZ^5.
\end{gathered}
\]

The ordering of the arrows is part of the result. The expressions are formed before they are assigned a physical realization; the ledgers are composed before the signature is extracted; the signature is evaluated through the character only after the memory has been composed. Thus, the composite domain is not reduced to the species included in a metrological table.

\section{Internal classification and compositional grammar}
\label{sec:alfabeto-composicional-rev6}

Let
\[
 \mathcal C_{81}=I_9\times I_9,
 \qquad I_9=\{1,\ldots,9\},
\]
the cell atlas, and let it be
\[
 \rho(r,c)=(10-r,10-c)
\]
its central inversion. We denote by
\(\mathcal L,\mathcal N,\mathcal D,\mathcal U\) the subsets of
charged leptons, neutrinos, down-type quarks, and up-type quarks,
respectively, and set
\[
 \mathcal Q=\mathcal D\sqcup\mathcal U.
\]
Each cell has a combinatorial level
\(\operatorname{gen}(x)\). In the quark sector the color
residue is also preserved
\[
 \kappa(x)=r(x)-c(x)\pmod3.
\]
We write
\[
 \mathcal Q_a=\{q\in\mathcal Q:\kappa(q)=a\},
 \qquad a\in\mathbb Z/3\mathbb Z.
\]
By \cref{prop:color-atlas},
\[
 |\mathcal Q_0|=|\mathcal Q_1|=|\mathcal Q_2|=12.
\]

\begin{definition}[Dual formal and cell identity]
\label{def:dual-identidad-composicional}
For an oriented cell section \(x\), its formal dual \(x^\vee\) inverts
the orientation of charge and color, and has support \(\rho(x)\). The
notation preserves the distinction
\[
 x^\vee\ne \rho(x)
\]
when \(\rho(x)\) is regarded as an unoriented cell. For each
\(x\in\mathcal C_{81}\), we denote by
\(\operatorname{id}_x:x\to x\) the identity arrow. The set of
identities is
\[
 \mathfrak I
 :=\{\operatorname{id}_x:x\in\mathcal C_{81}\}.
\]
\end{definition}

The identity \(\operatorname{id}_x\) is a neutral arrow of the compositional
system.
It is not, by that condition alone, a realization of a neutral boson. Likewise,
the formal dual incorporates orientation and cannot be replaced by a
second bare cell.

\begin{definition}[Mesonic and Baryonic Expressions]
\label{def:expresiones-mesonicas-barionicas}
The set of colour-neutral mesonic expressions is
\[
 \mathfrak M
 :=
 \left\{
   M(q,q')=q\otimes{(q')}^\vee:
   q,q'\in\mathcal Q,\ \kappa(q)=\kappa(q')
 \right\}.
\]
The pairs are ordered. Their total residue is
\(\kappa(q)-\kappa(q')=0\).

The set of baryonic expressions with fixed color positions is
\[
 \mathfrak B
 :=
 \mathcal Q_0\times\mathcal Q_1\times\mathcal Q_2,
\]
and we write
\[
 B(q_0,q_1,q_2)=q_0\otimes q_1\otimes q_2.
\]
Its residual neutrality is the identity.
\[
 0+1+2=0\pmod3.
\]
\end{definition}

These expressions are typed syntactic objects. In particular, no
quotient has yet been taken by permutations, spin, flavor, statistics,
binding energy, or dynamical equivalence.

\begin{definition}[Compatible loaded arrows]
\label{def:flechas-cargadas-compatibles}
The oriented leptonic arrows are
\[
 \mathfrak A_\ell
 :=
 \left\{
 (x,y)\in
   (\mathcal L\times\mathcal N)
   \sqcup
   (\mathcal N\times\mathcal L):
 \operatorname{gen}(x)=\operatorname{gen}(y)
 \right\}.
\]
Quark-oriented arrows are
\[
 \mathfrak A_q
 :=
 \left\{
 (x,y)\in
   (\mathcal U\times\mathcal D)
   \sqcup
   (\mathcal D\times\mathcal U):
 \begin{array}{l}
 \operatorname{gen}(x)=\operatorname{gen}(y),\\[-2pt]
 \kappa(x)=\kappa(y)
 \end{array}
 \right\}.
\]
We set
\[
 \mathfrak A_{\rm ch}:=\mathfrak A_\ell\sqcup\mathfrak A_q.
\]
\end{definition}

Level equality types the transition in both sectors; colour equality is
additionally required in the quark sector. Orientation distinguishes \(x\to y\)
from \(y\to x\).

\begin{definition}[Finite system of typed compositions]
\label{def:gramatica-celular-compuesta}
The finite compositional system is the typed disjoint union
\[
 \mathfrak G_{\rm comp}
 :=
 \mathcal C_{81}
 \sqcup\mathfrak I
 \sqcup\mathfrak M
 \sqcup\mathfrak B
 \sqcup\mathfrak A_{\rm ch}.
\]
The \(81\) cells are the primitive generators. The \(81\) identities are
arrows on those same generators, not a second atlas.
\end{definition}

\section{Combinatorial census of route compositions independent of metrological data}
\label{sec:censo-composicional-rev6}

\begin{theorem}[Census of the finite system of compositions]
\label{thm:censo-gramatica-composicional-rev6}
The system of the \cref{def:gramatica-celular-compuesta} contains
\[
 \boxed{\begin{gathered}
        81\ \text{identities},\qquad
        432\ \text{mesonic expressions},\\
        1728\ \text{baryonic expressions},\qquad
        372\ \text{charged arrows}
 \end{gathered}}
\]
The charged arrows decompose exactly as
\[
 372=320_{\text{leptonic}}+52_{\text{quark}}.
\]
The census depends only on the atlas, its inversion, the levels, and the
color residue. It uses neither masses, units, nor published signatures.
The four cardinalities count internal compositional expressions or admissible arrows
of the system; they do not count distinct physical species, physical particles,
or masses.
\end{theorem}

\begin{proof}
The assignment \(x\mapsto\operatorname{id}_x\) is bijective, so
\[
 |\mathfrak I|=|\mathcal C_{81}|=81.
\]
For mesons, the neutrality condition permits an ordered pair to be chosen
independently within each colour class. Since each
class contains twelve cells,
\[
 |\mathfrak M|
 =\sum_{a\in\mathbb Z/3\mathbb Z}|\mathcal Q_a|^2
 =3\cdot12^2
 =432.
\]
For baryons, the positions \(0,1,2\) are fixed, and therefore
\[
 |\mathfrak B|
 =|\mathcal Q_0||\mathcal Q_1||\mathcal Q_2|
 =12^3
 =1728.
\]

The charged lepton families have sizes
\[
 |\mathcal L_{G0}|=1,\qquad
 |\mathcal L_{G2}|=8,\qquad
 |\mathcal L_{G4}|=16,
\]
while the neutrino families have sizes
\[
 |\mathcal N_{G1}|=2,\quad
 |\mathcal N_{G2}|=4,\quad
 |\mathcal N_{G3}|=6,\quad
 |\mathcal N_{G4}|=8.
\]
Only \(G2\) and \(G4\) are common levels. Since both
orientations are counted,
\[
 |\mathfrak A_\ell|
 =2(8\cdot4+16\cdot8)
 =320.
\]

For the quark arrows, in the colour order \((0,1,2)\), the
up-type and down-type multiplicities at the common levels are
\[
\begin{array}{ccc}
 &\mathcal U&\mathcal D\\
 G1&(2,1,1)&(0,1,1)\\
 G3&(4,4,4)&(2,2,2).
\end{array}
\]
The number of pairs of an orientation and the same color is, therefore,
\[
 (2\cdot0+1\cdot1+1\cdot1)
 +(4\cdot2+4\cdot2+4\cdot2)
 =2+24=26.
\]
By including the opposite orientation one obtains
\[
 |\mathfrak A_q|=2\cdot26=52.
\]
Finally,
\[
 |\mathfrak A_{\rm ch}|=320+52=372.
\]
None of these rules contains an objective value or a mass vector.
\end{proof}

\begin{corollary}[Compositional constructor without whitespaces]
\label{cor:constructor-composicional-sin-blancos}
The constructor that jointly enumerates identities, mesonic expressions,
baryonic expressions, and charged arrows factorizes through the typed fields
\[
 (\text{sector},\text{family},\text{level},\kappa,\rho)
\]
of the atlas. With the preceding notation, its four codomains are
\(\mathfrak I,\mathfrak M,\mathfrak B,\mathfrak A_{\rm ch}\).
In particular, the census can be fixed before any signature table is opened.
\end{corollary}

\begin{proof}
The four definitions and the calculation of
\cref{thm:censo-gramatica-composicional-rev6} use only those fields.
The table of signatures does not belong to the domain of any of the rules.
\end{proof}

Target independence does not mean that the expressions have already acquired
a physical name. It means something prior and more precise: there exists a
compositional domain constructed without using as a criterion the quantity that will later be
evaluated.

\section{Memory transport from the composite route to the integral signature}
\label{sec:memoria-composicion-especies-rev6}

Let \(\gamma\in\Gamma_{108}^{\rm enr}\) be an enriched presentation of an
extension. Its read register contains
\[
 \mathcal E(\gamma)
 =
 (n_A,e_0,\ldots,e_{11},T_\zeta,C_\zeta),
\]
in addition to the phase, sheet, orientation, and edge data that type a composition.
The twelve events are not replaceable by their total before composition. For
\(0\le i<6\), we define
\[
 p_i=e_i+e_{i+6},
 \qquad
 a_i=e_i-e_{i+6},
\]
\[
 s_C=\sum_{i=0}^{5}p_i,
 \qquad
 M_i=p_i-p_5\quad(0\le i<5).
\]
With \(M_5=(M_0,\ldots,M_4)\) and
\(A_6=(a_0,\ldots,a_5)\), the reconstruction is
\[
 p_5=\frac{s_C-\sum_{i=0}^{4}M_i}{6},
 \qquad
 p_i=p_5+M_i,
\]
\[
 e_i=\frac{p_i+a_i}{2},
 \qquad
 e_{i+6}=\frac{p_i-a_i}{2}.
\]
By the \cref{thm:memoria-1-5-6}, on the integral image the
reversible decomposition is obtained
\[
 \boxed{(s_C,M_5,A_6),\qquad 1+5+6=12.}
\]
The first block preserves the visible total; the second, the relative form
within the symmetric sheet; the third, the memory between the two semicrowns.

The strong torsional coordinate also requires the cocycle condition at the step level:
\[
 k_{\Delta_4}=T_\zeta-2C_\zeta.
\]
It is not a function of the mere pattern of null events. This difference will be
essential when interpreting the non-descent rhombus.

\begin{definition}[Given oriented composition data]
\label{def:dato-empalme-orientado-rev6}
Let \(\Gamma,\Lambda\) be two enriched ledgers. A composition datum is a
pair \((g,B)\), where \(g\) aligns phase, sheet, orientation, and position at the
interface, and \(B\) is the common subtrace, with its contribution to the action, its
twelve events, its torsional predecessors, and its corona. On the domain in
which these data are compatible, one defines
\[
 \boxed{
 {(n,e)}_\Gamma\star_{(g,B)}{(n,e)}_\Lambda
 =
 \bigl(
 n_\Gamma+n_\Lambda-n_B,\,
 e_\Gamma+g e_\Lambda-e_B
 \bigr).}
\]
The torsional word and the corona are composed at the same level, before
computing \(T_\zeta,C_\zeta\) and \(k_{\Delta_4}\).
\end{definition}

The formula is an inclusion--exclusion of histories over an identified
subtrace. It is demonstrated with an explicit initial structure: the
records, the transport \(g\), the boundary \(B\), the orientation, and the word
of predecessors are part of the input. The compositional system does not
by itself choose any of those data.

\begin{definition}[Posterior extractor after the composition]
\label{def:extractor-posterior-empalme-rev6}
For a complete enriched record one defines
\[
 L_{\rm tick}(\gamma)
 =
 \bigl(
 n_A,s_C,k_{\Delta_4},
 \nu_{120}^{\rm ev},\nu_{270}
 \bigr)\in\ZZ^5,
\]
where, if \(\tau_m=\operatorname{sgn}(e_m)\) and
\(I_a=\{a,a+4,a+8\}\pmod {12}\),
\[
 \nu_{120}^{\rm ev}
 =
 \sum_{a=0}^{3}
 \operatorname{sgn}\left(\sum_{m\in I_a}e_m\right),
 \qquad
 \nu_{270}
 =
 \sum_{m=0}^{11}\tau_m\tau_{m+3}.
\]
The signature of a composition is defined by
\[
 \operatorname{Sig}(\Gamma,\Lambda;g,B)
 :=
 L_{\rm tick}
 \bigl(\Gamma\star_{(g,B)}\Lambda\bigr).
\]
\end{definition}

The signs and autocorrelations are applied to the already composed record.
Therefore, the definition is not equivalent to summing the signatures of the constituents.

\section{Commutative diagram of composition and fiber preservation}
\label{sec:rombo-no-descenso-composicional-rev6}

The historical catalog B1 employs an event projection that is convenient to type separately. Let
\[
 \Sigma_{12}=(+,+,+,-,-,-,+,+,+,-,-,-),
 \qquad \tau_m=\operatorname{sgn}(e_m),
\]
and let us define the null event proxy
\[
 K_0(e)
 =
 \sum_{\tau_m=0}\Sigma_{12}(m).
\]
The exact projection is
\[
 F_0(n_A,e)
 =
 \left(
 n_A,\,
 \sum_m e_m,\,
 K_0(e),\,
 \sum_{a=0}^{3}
   \operatorname{sgn}\sum_{m\in I_a}e_m,\,
 \sum_m\tau_m\tau_{m+3}
 \right)\in\ZZ^5.
\]
For two aligned records, we write
\[
 (n,e)\boxplus(n',e')=(n+n',e+e').
\]
This operation is the case with no common subtrace and with identity alignment of the
event-composition case.

\begin{proposition}[Obstruction to the descent of the projection B1]
\label{prop:no-descenso-proyeccion-b1-rev6}
There is no binary operation
\[
 \mu:\ZZ^5\times\ZZ^5\longrightarrow\ZZ^5
\]
such that
\[
 F_0(\Gamma\boxplus\Lambda)
 =
 \mu\bigl(F_0(\Gamma),F_0(\Lambda)\bigr)
\]
for all records B1.
\end{proposition}

\begin{proof}
Consider the routes, written in the chart
\((r,c,h_0,v_0)\),
\[
 a=(1,3,-1,+1),
 \qquad
 b=(6,7,-1,+1),
 \qquad
 c=(1,2,-1,-1).
\]
The first two have distinct event records:
\[
\begin{aligned}
 e(a)&=(5,4,4,-4,-4,-6,\;5,4,4,-4,-4,-6),\\
 e(b)&=(4,5,4,-4,-6,-4,\;4,5,4,-4,-6,-4),
\end{aligned}
\]
but publish the same five-tuple:
\[
 F_0(a)=F_0(b)=(8,-2,0,0,-12).
\]
The test route has
\[
 e(c)=(1,4,1,-2,-4,-3,\;1,4,1,-2,-4,-3)
\]
and
\[
 F_0(c)=(28,-6,0,-4,-12).
\]
By first summing the events and then reevaluating, one obtains
\[
 F_0(a\boxplus c)=(36,-8,0,0,-12),
\]
\[
 F_0(b\boxplus c)=(36,-8,0,-2,-12).
\]
The two pairs of inputs of the putative \(\mu\) are equal, but the
outputs differ in \(\nu_{120}^{\rm ev}\). This contradicts the existence of
\(\mu\).
\end{proof}

The exhaustive census of the same domain contains \(324\) routes, \(53\)
proxy five-tuples, and \(198\) distinct event states. Forty-two
fibers of the projection contain more than one enriched state; forty of
them are separated by some test route. The preceding diamond is the shortest
canonical witness to that loss of information.

\begin{remark}[Typed separation between the representative \(K_0\) and strong torsion]
\label{rem:proxy-no-torsion-fuerte-rev6}
The third coordinate of \(F_0\) is \(K_0(e)\), which in B1 takes values in
\(\{-4,-2,0,2,4\}\). It is not identified with
\(k_{\Delta_4}=T_\zeta-2C_\zeta\). Indeed, there exist routes with the same
\(n_A\) and the same vector of twelve events to which the torsional reader
assigns distinct values because it preserves the predecessors of each iteration
and the crown.

The \cref{prop:no-descenso-proyeccion-b1-rev6} proves exactly that the
five-coordinate B1 projection does not transport the superposition of
records. It does not assert a universal impossibility for every future structure
on \(\ZZ^5\), nor does it authorize replacing \(K_0\) with
\(k_{\Delta_4}\). Under the strong reader, the output is defined through the
\cref{def:extractor-posterior-empalme-rev6}: the enriched word is first
composed and its five coordinates are then re-extracted.
\end{remark}

\section{Theorem of independent data record composition in metrology}
\label{sec:teorema-composicional-sin-blancos-rev6}

\begin{theorem}[Compositional closure without blanks]
\label{thm:cierre-composicional-sin-blancos-rev6}
Once the atlas \(\mathcal C_{81}\), formal duality, the color residue,
and the enriched record of the \cref{chap:extension-ruta-caracter} have been fixed, the following
assertions hold.
\begin{enumerate}
 \item There exists the finite system of typed compositions
 \(\mathfrak G_{\rm comp}\) with \(81\) identities, \(432\) mesonic expressions,
 \(1728\) baryonic expressions, and \(372\) charged arrows,
 decomposed as \(320+52\).

 \item The system is constructed without target masses. Its output is a
 typed expression, orbit, or multisection, not a mass or an individualized
 physical label.

 \item Given an expression and constituent records with an
 explicit tree of compatible compositions \((g,B)\), the operation produces an
 enriched record. Its signature is obtained by the application
 \[
 \boxed{
 \bigl(
 \text{expression},
 \text{records},
 \text{oriented compositions}
 \bigr)
 \xrightarrow{\ \mathfrak C_{12}\ }
 \text{composite record}
 \xrightarrow{\ L_{\rm tick}\ }
 \ZZ^5.}
 \]
 Here \(\mathfrak C_{12}\) denotes the iteration of
 \(\star_{(g,B)}\) with the parenthesization and boundaries included in the data.

 \item The B1 projection onto five integers does not by itself constitute a
 compositional space: the diagram with \(\boxplus\) does not descend to
 \(\ZZ^5\). Therefore, replacing each record by its five-tuple before
 composition loses information needed to reevaluate at least
 \(\nu_{120}^{\rm ev}\).
\end{enumerate}
\end{theorem}

\begin{proof}
The first point is \cref{thm:censo-gramatica-composicional-rev6}. The second is \cref{cor:constructor-composicional-sin-blancos}. For the third, each edge of the composition tree is evaluated through \cref{def:dato-empalme-orientado-rev6}; because the boundaries and transports form part of the domain, the iteration produces a register on which \(L_{\rm tick}\) is defined. No intermediate signature is used as a substitute for the register. The fourth point is \cref{prop:no-descenso-proyeccion-b1-rev6}.
\end{proof}

The correct formula for the composite signature is, therefore,
\[
 \operatorname{Sig}_{\rm comp}
 =
 L_{\rm tick}\circ\mathfrak C_{12}.
\]
It is not a sum of mass rows. Only afterward may the formal character
or its global specialization be applied:
\[
 \text{composite record}
 \xrightarrow{L_{\rm tick}}\ZZ^5
 \xrightarrow{\chi_{\rm form}}\operatorname{Fun}(\mathcal P,\RR_{>0})
 \longrightarrow
 \text{dimensional chart}.
\]

\section{Genealogical reconstruction of mixing, generation, and persistence of flavor}
\label{sec:significado-genealogico-composicion-rev6}

The preceding theorem completes a stratum that a list of masses cannot
display. Primitive cells do not merely receive names; they generate
identities, dual pairs, color triples, and oriented arrows through
internal rules. Composite objects thus appear within the domain before
any metrological comparison.

Genealogy is not reduced to the syntactic tree either. Two expressions with the
same number of constituents may be composed in different ways because
the phase, sheet, orientation, and common subtrace form part of the datum.
The register preserves those differences until the signs,
correlations, and torsion have acted. Mass, when a dimensional
realization exists, will be a subsequent reading of that composite memory.

This order explains why \(\ZZ^5\) is not the composition domain although
it is the domain of the character. The group \((\ZZ^5,+)\) exists, and the character is
multiplicative on it; what does not exist for the B1 projection is a law
that recovers from the five-tuples the superposition that has already forgotten
the placement of the twelve events. Linearity of the character does not restore the
information removed by the extractor.

\section{From the internal compositional system to physical realization}
\label{sec:delimitacion-composicion-especies-rev6}

The \cref{thm:cierre-composicional-sin-blancos-rev6} constructs the domain
prior to physical names: ordered expressions, duals, color triples, and
compatible arrows. Their cardinalities count internal possibilities of
composition. A physical species, by contrast, is a persistent representation
with specified flavor, color, spin, statistics, and channel. Conflating the two
levels would turn a census of expressions into a spurious catalogue of
particles.

\begin{definition}[Physical realization data]
\label{def:dato-realizacion-fisica-rev6}
A realization datum for a species \(X\) is the tuple
\[
 \mathfrak D_X=
 \bigl(
 \mathfrak t_X,\mathfrak s_X,
 (\mathcal L_{X_i})_{i=1}^{r},
 \mathcal T_X,(g_X,B_X)
 \bigr),
\]
where:
\begin{enumerate}
 \item \(\mathfrak t_X\) fixes the physical type ---elementary, null, mesonic,
 baryonic, multiquark, or topological--- and its representation;
 \item \(\mathfrak s_X\) is the multisection or realized internal expression;
 \item \(\mathcal L_{X_i}\) are the enriched ledgers of the
 constituents;
 \item \(\mathcal T_X\) fixes flavor, color, spin, statistics, and
 symmetrization;
 \item \((g_X,B_X)\) fixes the composition tree, its transports, and its
 common subtraces.
\end{enumerate}
The tuple is declared without consulting the mass or the width, which will be evaluated later.
\end{definition}

\begin{proposition}[The realization changes the register, not the law]
\label{prop:realizacion-compuesta-ley-unica-rev6}
Every datum \(\mathfrak D_X\) determines a composite record
\[
 \mathcal L_X
 =\mathfrak C_{12}
   \bigl((\mathcal L_{X_i})_{i=1}^{r};
         \mathcal T_X,g_X,B_X\bigr),
\]
a signature \(v_X=L_{\rm tick}(\mathcal L_X)\) and the universal central character
\[
 \chi_0(v_X-v_e).
\]
If the genealogy also provides a typed spectral refinement
\(\eta_X\in\mathcal E_{\mathfrak S}\), the complete evaluation is
\[
 m_X
 =m_e^{\rm HMT}\,
   \chi_{\rm full}(v_X-v_e,\eta_X).
\]
The central character and the refinement belong to distinct levels; the
completeness of the tower does not replace the prospective map that produces
\(\eta_X\).
If \(X\) is unstable, its width belongs to the temporal reader
\(\Gamma_X=\hbar_{\rm int}/\tau_X^{\rm life}\), not to a sixth coordinate of
the mass signature.
\end{proposition}

\begin{proof}
The realization datum contains exactly the entries of
\(\mathfrak C_{12}\). The
\cref{thm:cierre-composicional-sin-blancos-rev6} produces the enriched
record, and the \cref{def:extractor-posterior-empalme-rev6} produces its
signature. The central character and its spectral lift are those of
\cref{eq:caracter-masa-global-rev7,eq:ley-masa-todos-registros-rev6}. The
width uses a temporal line and therefore has a different domain.
\end{proof}

The level separation is thus fixed:
\[
\boxed{
\begin{gathered}
\text{internal compositional system}
\longrightarrow\text{realization datum}
\longrightarrow\text{composite record}\\
\longrightarrow\text{signature}
\longrightarrow\text{central character}
\longrightarrow\text{spectral refinement}
\longrightarrow\text{physical comparison}.
\end{gathered}}
\]
The \(432\) mesonic expressions, the \(1728\) baryonic expressions, and the \(372\)
charged arrows occupy only the first slot and do not constitute an
inventory of particles. The Standard Model preserves its typed set of
species, and each species occupies the chain through its own realization
datum. The causal rule is unique across all sectors:
\[
 \boxed{\text{compose extensions and records;\quad read the signature afterwards}.}
\]
